\section*{উৎস ও পাঠের পরিধি}
\addcontentsline{toc}{section}{উৎস ও পাঠের পরিধি}
এই অধ্যায়ের ব্যাখ্যা ও অনুশীলন নিজস্ব ভাষায় লেখা। মাহবুবুর রহমানের সরবরাহকৃত 2026 সালের বইয়ের Chapter 1 (মুদ্রিত পৃষ্ঠা 1--98; PDF পৃষ্ঠা 7--104) ধারণা ও প্রশ্নের পরিধি মিলিয়ে দেখা হয়েছে।
\begin{itemize}
\item AI: Russell ও Norvig-এর \href{https://aima.cs.berkeley.edu/contents.html}{Artificial Intelligence: A Modern Approach-এর প্রকাশ্য সূচিপত্র}; \href{https://www-formal.stanford.edu/jmc/history/dartmouth/dartmouth.html}{Dartmouth-এর মূল গবেষণা-প্রস্তাব}। বাণিজ্যিক বইটির পূর্ণ পাঠ পর্যালোচনার দাবি করা হচ্ছে না।
\item চিকিৎসা: \href{https://www.cancer.gov/about-cancer/treatment/types/surgery/cryosurgery}{NCI Cryosurgery}; জৈব প্রযুক্তি: \href{https://www.genome.gov/genetics-glossary/Genetic-Engineering}{NHGRI Genetic Engineering}, \href{https://www.genome.gov/about-genomics/fact-sheets/Deoxyribonucleic-Acid-Fact-Sheet}{DNA}, \href{https://www.ncbi.nlm.nih.gov/genbank/genbank/}{NCBI GenBank ও BLAST}।
\item Biometrics: \href{https://nvlpubs.nist.gov/nistpubs/SpecialPublications/NIST.SP.800-76-2.pdf}{NIST-এর Biometric Data Specification}; Nano: \href{https://www.nano.gov/about-nanotechnology/}{NNCO} এবং \href{https://feynman.com/science/nanotechnology/}{Feynman-এর বক্তৃতা}।
\item Navigation: \href{https://www.gps.gov/systems/gps/index.php}{GPS.gov}; RFID: \href{https://www.gs1.org/standards/rfid}{GS1}; নৈতিকতা: \href{https://computerethics.institute/publications/ten-commandments-of-computer-ethics/}{Computer Ethics Institute}।
\item Robotics-এর দেশতুলনায় \href{https://ifr.org/ifr-press-releases/news/global-robot-density-in-factories-doubled-in-seven-years}{IFR-এর 2023 Robot Density} ও \href{https://ifr.org/ifr-press-releases/news/global-robot-demand-in-factories-doubles-over-10-years}{2024 Installation Data} উদাহরণ হিসেবে ব্যবহৃত। এগুলো স্থায়ী Ranking নয়।
\item প্রশ্নের ধরন: \href{https://ictshikhi.com/mcq/hsc/information-and-communication-technology-world-and-bangladesh-perspective}{ICT Shikhi-এর প্রকাশ্যে দৃশ্যমান প্রথম 20টি MCQ} এবং \href{https://ictshikhi.com/cq/hsc/information-and-communication-technology-world-and-bangladesh-perspective}{8টি বিনামূল্যে দৃশ্যমান CQ}। এগুলো তৃতীয় পক্ষের পুনর্মুদ্রণ; বোর্ডের মূল প্রশ্নপত্র হিসেবে স্বাধীনভাবে সত্যতা নিশ্চিত করা হয়নি।
\end{itemize}
\begin{itemize}
\item আগের প্রশ্নের অতিরিক্ত যাচাই: \href{https://newtonpublications.com/ebook/hsc_tp_me/common/ICT.pdf}{Newton-এর প্রকাশ্য নমুনায় 2018--2019 সালের প্রশ্নপত্র}, TeachingBD-এর 2016--2017 সালের আট বোর্ডের স্ক্যান এবং \href{https://onetimeschool.com/hsc-exam/hsc-information-and-communication-technology-multiple-choice-question-solution-2024/4323/}{2024 সালের নয় বোর্ডের MCQ স্ক্যান}। পুনর্মুদ্রিত উত্তরমালা সরাসরি গ্রহণ না করে Concept ও প্রামাণ্য উৎসের সঙ্গে মিলিয়ে দেখা হয়েছে।
\item ইতিহাস ও বিশ্লেষণ: \href{https://www.genome.gov/25520302/online-education-kit-1972-first-recombinant-dna}{NHGRI Recombinant DNA ইতিহাস}; \href{https://arxiv.org/abs/2111.11832}{Hesper ও Hogeweg-এর 1970 সালের মূল লেখার অনুবাদ}; \href{https://pubmed.ncbi.nlm.nih.gov/11552348/}{Bioinformatics-এর পদ্ধতি নিয়ে গবেষণা}; \href{https://www.genome.gov/genetics-glossary/Gene-Mapping}{NHGRI Gene Mapping}।
\item পরিভাষা ও বাংলাদেশের গবেষণা: \href{https://www.catb.org/~esr/sf-words/glossary.html}{Science Fiction Lexicography-তে Genetic Engineering শব্দের সাহিত্যিক ব্যবহার}; \href{https://www.nature.com/articles/nplants2016223}{পাটের দুই প্রজাতির Genome নিয়ে মূল গবেষণা}; \href{https://cdn.most.gov.bd/image/64accaccd1f27}{বিজ্ঞান ও প্রযুক্তি মন্ত্রণালয়ের প্রকাশিত গবেষণা-পর্যালোচনায় 2010 সালের পাটের Genome উদ্যোগ}।
\item পরিচিত পণ্যের উদাহরণ: \href{https://www.nano.gov/sites/default/files/ScavengerHunt.pdf}{NNCO-এর Tennis Ball ও Nano Coating ব্যাখ্যা}; \href{https://learn.microsoft.com/en-us/copilot/glossary}{Microsoft-এর Copilot পরিভাষা ও Natural Language Assistant ব্যাখ্যা}। Product-এর নাম সময়ভেদে বদলাতে পারে; উদাহরণ দিয়ে মূল প্রযুক্তি চেনা শেখানো হয়েছে।
\item প্রয়োগ: \href{https://www.nano.gov/about-nanotechnology/applications-nanotechnology/}{NNCO Nano Applications}, \href{https://www.universal-robots.com/media/1828411/ur-product-factsheets-ur10e-06_2023-en-version-5-1-4-3.pdf}{Robot Arm-এর নির্মাতার তথ্য}, \href{https://www.bostonscientific.com/content/dam/elabeling/pi/visual-ice/51696695-01A_VISUAL-ICE_Manual_ML_s.pdf}{Cryoablation System-এর নির্মাতার নির্দেশিকা}। চিকিৎসার উদাহরণ প্রযুক্তির কাজ বোঝায়।
\item যোগাযোগ ও উৎপাদন: \href{https://news.njit.edu/murray-turoff-pioneer-instant-messaging}{NJIT-এর Murray Turoff বিবরণ}, \href{https://home.cern/science/computing/the-birth-of-the-web/}{CERN-এর Web ইতিহাস}, \href{https://www.ibm.com/think/topics/industry-4-0}{IBM-এর Industry 4.0 ব্যাখ্যা}; চিকিৎসার চিত্র: \href{https://www.radiologyinfo.org/en/info/pediatric-mri}{MRI} ও \href{https://www.radiologyinfo.org/en/info/imrt}{Computer-নিয়ন্ত্রিত Radiotherapy}।
\end{itemize}
সব বছরের সব Set-এর মূল প্রশ্নপত্র পাওয়া যায়নি। কোন বছর ও বোর্ডের পূর্ণ বা আংশিক উৎস মিলেছে, প্রশ্নের কোন ধারণা কোথায় ব্যাখ্যা করা আছে এবং কোথায় যাচাই বাকি, তা অধ্যায়ের সঙ্গে রাখা \texttt{CHAPTER01\_BOARD\_AUDIT.md} ও তার Evidence তালিকায় নথিভুক্ত।
পরিভাষা, সংজ্ঞা, কাজের ধাপ, পার্থক্য ও উদ্দীপকের সূত্র মূল টপিকে রাখা হয়েছে। নতুন অনুশীলনগুলো নিজস্ব প্রশ্ন; নির্দিষ্ট বোর্ডের হুবহু প্রশ্ন নয়। অতিরিক্ত উদাহরণ ধারণা বোঝার জন্য, সব নাম বা Product মুখস্থ করার নির্দেশ নয়। যাচাইয়ের তারিখ: 5 October 2026।

\subsubsection{আধুনিক AI অংশের অতিরিক্ত উৎস}
এই অংশের Model পরিচিতিতে নির্মাতাদের প্রকাশ্য নথি, এবং ধারণা বোঝাতে মূল গবেষণা ব্যবহার করা হয়েছে। নির্দিষ্ট Version-এর দাম, Context Limit বা Ranking স্থায়ী তথ্য হিসেবে দেওয়া হয়নি। উদাহরণ, সমস্যা ও সমাধানের বর্ণনা নিজস্ব ভাষায় লেখা।
\begin{compactpoints}
\item \href{https://arxiv.org/abs/1706.03762}{Vaswani ও সহকর্মী: Attention Is All You Need (2017)}; \href{https://arxiv.org/abs/2005.14165}{Brown ও সহকর্মী: Language Models are Few-Shot Learners (2020)}---Transformer ও Language Model-এর বিকাশ।
\item \href{https://arxiv.org/abs/2005.11401}{Lewis ও সহকর্মী: Retrieval-Augmented Generation (2020)}---Retrieval দিয়ে প্রাসঙ্গিক তথ্য যুক্ত করার ধারণা।
\item \href{https://arxiv.org/abs/2112.10752}{Rombach ও সহকর্মী: Latent Diffusion}; \href{https://arxiv.org/abs/2212.04356}{Radford ও সহকর্মী: Speech Recognition/Whisper}---Image ও Speech Model-এর ভিন্ন কাজ।
\item \href{https://arxiv.org/abs/2210.03629}{Yao ও সহকর্মী: ReAct}; \href{https://www.anthropic.com/engineering/building-effective-agents}{Anthropic-এর Agent/Workflow ব্যাখ্যা}; \href{https://developers.openai.com/api/docs/guides/tools}{OpenAI Tool নথি}---Tool Action, Observation ও Agent-এর Loop।
\item \href{https://arxiv.org/abs/2304.11771}{Brynjolfsson, Li ও Raymond: Generative AI at Work}; \href{https://deepmind.google/science/alphafold/}{AlphaFold}; \href{https://www.unesco.org/en/articles/guidance-generative-ai-education-and-research}{UNESCO-এর শিক্ষা ও গবেষণায় GenAI নির্দেশনা}---অবদান ও তার সীমা।
\end{compactpoints}
Model Family-র প্রামাণ্য Link সংশ্লিষ্ট সারণির নিচে আছে। ভবিষ্যৎ উন্নয়নের আলোচনা বর্তমান গবেষণার দিক থেকে শর্তযুক্ত সম্ভাবনার ব্যাখ্যা; কোনো নিশ্চিত সময়সীমা বা সম্পূর্ণ স্বায়ত্তশাসনের দাবি নয়।

